\section{Certified approximate light cones}
\label{sec:approx-lightcone}

The exact light-cone theorem requires literal zero phase couplings.  Real LD,
annotation, and feature graphs are often only approximately sparse.  This
section replaces exact support by a computable phasewise envelope.

Let the phase projectors be $P_1,\ldots,P_q$, and let $X$ be a vector or block
vector.  Define its phase-norm profile
\[
\rho(X)_g=\lVert P_gX\rVert_F,
\]
and define the nonnegative phase-transfer matrix
\[
\mathcal B_{hg}=\lVert P_hHP_g\rVert_2.
\]
The matrix $\mathcal B$ records the strength of every one-step phase transfer
without assuming that weak couplings vanish.

\begin{theorem}[Weighted phase-transfer envelope]
\label{thm:weighted-transfer}
For every block vector $X$,
\[
\boxed{
\rho(HX)\leq \mathcal B\rho(X)
}
\]
componentwise.  Consequently, for every integer $k\geq0$,
\[
\boxed{
\rho(H^kX)\leq \mathcal B^k\rho(X).
}
\]
\end{theorem}

\begin{proof}
Using $\sum_gP_g=I$ and submultiplicativity,
\[
\begin{aligned}
\lVert P_hHX\rVert_F
&=\left\lVert\sum_gP_hHP_gX\right\rVert_F\\
&\leq\sum_g\lVert P_hHP_g\rVert_2\lVert P_gX\rVert_F\\
&=(\mathcal B\rho(X))_h.
\end{aligned}
\]
Iteration gives the power bound.
\end{proof}

Now choose a symmetric set of retained phase blocks and define
\[
H_0=\sum_{(h,g)\in\mathcal E_0}P_hHP_g.
\]
Let
\[
(\mathcal B_0)_{hg}=
\begin{cases}
\mathcal B_{hg},&(h,g)\in\mathcal E_0,\\
0,&\text{otherwise},
\end{cases}
\qquad
\mathcal R=\mathcal B-\mathcal B_0.
\]
Symmetric pair retention preserves self-adjointness of $H_0$, although the
following error theorem does not require positive semidefiniteness of $H_0$.

\begin{theorem}[Thresholded-power certificate]
\label{thm:thresholded-power}
For every $k\geq1$,
\[
\boxed{
\rho\bigl((H^k-H_0^k)X\bigr)
\leq
(\mathcal B^k-\mathcal B_0^k)\rho(X).
}
\]
In particular,
\[
\boxed{
\lVert(H^k-H_0^k)X\rVert_F
\leq
\left\lVert(\mathcal B^k-\mathcal B_0^k)\rho(X)\right\rVert_2.
}
\]
\end{theorem}

\begin{proof}
Write $E=H-H_0$.  The telescoping identity gives
\[
H^k-H_0^k
=
\sum_{j=0}^{k-1}H^{k-1-j}EH_0^j.
\]
The block-norm matrix of $E$ is $\mathcal R$.  Applying
Theorem~\ref{thm:weighted-transfer} successively to each term gives
\[
\rho\bigl((H^k-H_0^k)X\bigr)
\leq
\sum_{j=0}^{k-1}
\mathcal B^{k-1-j}\mathcal R\mathcal B_0^j\rho(X).
\]
The matrix telescoping identity
\[
\mathcal B^k-\mathcal B_0^k
=
\sum_{j=0}^{k-1}
\mathcal B^{k-1-j}(\mathcal B-\mathcal B_0)\mathcal B_0^j
\]
yields the phasewise formula.  Orthogonality of the phase projectors gives
\[
\lVert Y\rVert_F^2=\sum_h\lVert P_hY\rVert_F^2,
\]
which proves the global bound.
\end{proof}

\begin{corollary}[Certified approximate light cone]
\label{cor:approx-lightcone}
Assume $X$ is supported on a set of start phases $S$.  If phase $h$ is not
reachable from $S$ by a retained path of length at most $k$, then
\[
P_hH_0^kX=0
\]
and
\[
\boxed{
\lVert P_hH^kX\rVert_F
\leq
\bigl[(\mathcal B^k-\mathcal B_0^k)\rho(X)\bigr]_h.
}
\]
Thus an exact zero outside the retained light cone is replaced by a certified
nonzero leakage budget.
\end{corollary}

\begin{corollary}[Scalar threshold law]
For any induced matrix norm, let
\[
M=\max\{\lVert\mathcal B\rVert,\lVert\mathcal B_0\rVert\}.
\]
Then
\[
\boxed{
\lVert\mathcal B^k-\mathcal B_0^k\rVert
\leq
kM^{k-1}\lVert\mathcal R\rVert.
}
\]
Hence weak discarded coupling produces a first-order error bound proportional
to the discarded transfer norm and to the Krylov depth.
\end{corollary}

\begin{theorem}[Krylov-polynomial certificate]
\label{thm:polynomial-certificate}
Let
\[
p(t)=\sum_{j=0}^d a_jt^j.
\]
Then
\[
\boxed{
\rho\bigl((p(H)-p(H_0))X\bigr)
\leq
\sum_{j=1}^d
|a_j|(\mathcal B^j-\mathcal B_0^j)\rho(X).
}
\]
\end{theorem}

\begin{proof}
Apply Theorem~\ref{thm:thresholded-power} to every nonconstant power and use
the triangle inequality.  The constant terms cancel.
\end{proof}

The polynomial theorem covers arbitrary vectors in a finite Krylov space,
including polynomial filters and restarted combinations, rather than only the
monomial $H^kX$.

\begin{corollary}[Certified phase pruning]
If a thresholded phase graph is chosen so that
\[
\left\lVert
(\mathcal B^k-\mathcal B_0^k)\rho(X)
\right\rVert_2
\leq\tau,
\]
then replacing $H^kX$ by $H_0^kX$ incurs Frobenius error at most $\tau$.
The same statement holds for a Krylov polynomial with the certificate in
Theorem~\ref{thm:polynomial-certificate}.
\end{corollary}

This turns graph thresholding into an error-budget decision.  A phase or
interface is not discarded because it merely appears small; it is discarded
only when the weighted sum of all paths using discarded couplings remains
inside the requested tolerance.

\begin{figure}[H]
\centering
\includegraphics[width=.84\linewidth]{figure_09_approximate_lightcone_certificate.png}
\caption{Weak long-range phase couplings replace the exact zero outside the retained light cone by a nonzero but certified envelope.  In this PSD test, the phasewise transfer-matrix certificate bounds the observed leakage at every depth from one through eight.}
\label{fig:approx-lightcone-certificate}
\end{figure}

\begin{figure}[H]
\centering
\includegraphics[width=.84\linewidth]{figure_10_threshold_tradeoff.png}
\caption{Certified phase-graph sparsification at Krylov depth five.  Retaining fewer off-diagonal phase blocks reduces communication and operator work while the transfer-matrix bound remains above the measured truncation error.  The bound is intentionally conservative because it ignores cancellation and singular-vector alignment.}
\label{fig:threshold-tradeoff}
\end{figure}

\subsection{Numerical audit}

The validation operator has twelve phases of width four.  A block-bidiagonal
factor generates a positive-definite local operator $H_0$, and weak symmetric
long-range blocks are then added to obtain $H$.  Both normalized operators
remain positive definite.  A three-column start block is supported in phase
zero.

For depths one through eight:
\begin{itemize}
\item every observed phase norm lay below $\mathcal B^k\rho(X)$;
\item every observed thresholding error lay below
$(\mathcal B^k-\mathcal B_0^k)\rho(X)$;
\item the largest total observed-error-to-certificate ratio was approximately
$0.669$.
\end{itemize}

A separate hard-threshold sweep retained between $22$ and $132$ directed
off-diagonal phase blocks.  Every measured depth-five error remained below its
certificate; the largest observed-error-to-certificate ratio was approximately
$0.376$.  Finally, a degree-four signed polynomial in $H$ satisfied the
phasewise polynomial certificate in all twelve phases, with maximum ratio
approximately $0.068$.

These computations validate the implementation and illustrate sharpness.  The
universal result is supplied by the preceding proofs, not by the finite test.
